\documentclass[12pt,a4paper]{article}
\usepackage[utf8]{inputenc}
\usepackage[T1]{fontenc}
\usepackage[polish]{babel}
\usepackage{geometry}
\usepackage{enumitem}
\usepackage{titlesec}
\usepackage{parskip}
\usepackage{fancyhdr}
\usepackage{xcolor}
\usepackage{mdframed}
\usepackage{microtype}
\usepackage{booktabs}
\usepackage{hyperref}
\usepackage{multirow}
\usepackage{array}
\usepackage{longtable}
\usepackage{tabularx}

\geometry{margin=2.5cm}
\hypersetup{colorlinks=true, linkcolor=black, urlcolor=blue}

\pagestyle{fancy}
\fancyhf{}
\rhead{SDLC-LLM Benchmark}
\lhead{Zadanie R2 -- Refinement (Hybrid)}
\rfoot{\thepage}

\titleformat{\section}{\large\bfseries}{}{0em}{}[\titlerule]
\titleformat{\subsection}{\normalsize\bfseries}{}{0em}{}
\titleformat{\subsubsection}{\normalsize\itshape}{}{0em}{}

\definecolor{metabg}{gray}{0.93}
\definecolor{promptbg}{RGB}{235,245,255}
\definecolor{score3}{RGB}{220,240,220}
\definecolor{score2}{RGB}{255,243,205}
\definecolor{score1}{RGB}{255,220,200}
\definecolor{score0}{RGB}{255,200,200}
\definecolor{tpcolor}{RGB}{220,240,220}
\definecolor{fpcolor}{RGB}{255,220,200}
\definecolor{fncolor}{RGB}{255,243,205}

\newcommand{\scorebox}[1]{%
  \ifnum#1=3\colorbox{score3}{\textbf{3}}%
  \else\ifnum#1=2\colorbox{score2}{\textbf{2}}%
  \else\ifnum#1=1\colorbox{score1}{\textbf{1}}%
  \else\colorbox{score0}{\textbf{0}}\fi\fi\fi
}
\newcommand{\tp}[1]{\colorbox{tpcolor}{\small #1}}
\newcommand{\fp}[1]{\colorbox{fpcolor}{\small #1}}
\newcommand{\fn}[1]{\colorbox{fncolor}{\small #1}}

\begin{document}

% ─── STRONA TYTUŁOWA ────────────────────────────────────────────────
\begin{center}
  {\LARGE\bfseries Requirements Artifact}\\[1.5em]
  {\large \textbf{Case study:} Appointment Booking System}\\[0.3em]
  {\large \textbf{Model:} \texttt{gpt-5.4}}\\[0.3em]
  {\large \textbf{Tryb:} \textbf{Hybrid} (LLM + Human-in-the-loop)}\\[0.3em]
  {\large \textbf{Wejście:} Gold ver2 (challenging -- celowe niejednoznaczności)}
\end{center}

\vspace{1em}\hrule\vspace{1.5em}

% ─── WPROWADZONE ZMIANY (HUMAN CHANGES) ─────────────────────────────
\section{Interwencja analityka (Human Changes)}

\begin{mdframed}[backgroundcolor=promptbg, linewidth=0pt, innerleftmargin=10pt, innerrightmargin=10pt, innertopmargin=8pt, innerbottommargin=8pt]
\small
\textbf{Changelog z przeglądu manualnego:}\\
FR3: zmodyfikowano ogólne warunki anulowania, dodając sztywny limit 'do 24 godzin przed zaplanowanym terminem'. 
FR4: zastąpiono stary wymóg definiowania ograniczeń nowym wymogiem o powrocie anulowanego terminu do puli dostępnych. 
FR5: doprecyzowano zakaz konfliktów poprzez jawne zablokowanie podwójnych rezerwacji. 
FR9: zastąpiono generyczny wymóg obsługi ról ścisłą definicją statusów wizyt (AVAILABLE, BOOKED, CANCELLED, COMPLETED). 
FR10: zastąpiono ogólny wymóg skalowalności kluczowym wymogiem biznesowym o limicie maksymalnie 3 aktywnych rezerwacji. 
NFR1: zastąpiono ogólną dostępność konkretnym wymogiem czasu odpowiedzi poniżej 1 sekundy. 
NFR4: zastąpiono ogólne zasady RBAC ścisłym wymogiem uwierzytelnienia. 
NFR5: zaktualizowano log audytowy (dawne NFR6) o znacznik czasu. 
NFR ogólne: usunięto NFR5, NFR7, NFR8 (dotyczące ogólnego bezpieczeństwa, konfiguracji UI i czytelności interfejsu), skracając listę z 8 do 5. 
US ogólne: doprecyzowano limit 24h w US3, usunięto US4, US7, US8 (dotyczące informacji o limitach w UI, dedykowanych widoków specjalisty oraz zarządzania rolami), przenumerowano listę z 10 do 7 punktów. 
AC ogólne: usunięto rozbudowane kryteria UI i poboczne (stare AC1, AC2, AC7, AC8, AC9, AC10, AC11, AC12), zmodyfikowano kryteria anulowania wprowadzając logikę 24h (nowe AC3, AC4), dodano nowe AC5 weryfikujące limit 3 rezerwacji. Zredukowano liczbę AC z 14 do 7, pozostawiając tylko kluczowe przepływy rezerwacji, współbieżności i konfiguracji administratora.
\end{mdframed}

% ─── CEL SYSTEMU ────────────────────────────────────────────────────
\section{Cel systemu}

System umożliwia użytkownikom rezerwację, przeglądanie i anulowanie wizyt u specjalistów oraz wspiera specjalistów w zarządzaniu grafikiem. Celem systemu jest bezpieczna i skalowalna obsługa terminów, rezerwacji i zmian harmonogramu przy uwzględnieniu ról użytkowników, ograniczeń czasowych oraz wyjątków od konfliktów rezerwacji.

% ─── ROLE ───────────────────────────────────────────────────────────
\section{Role}

\begin{itemize}[leftmargin=1.5em]
  \item \textbf{Użytkownik} -- przeglądanie dostępnych terminów, dokonywanie rezerwacji wizyt, anulowanie własnych rezerwacji w dozwolonych warunkach.
  \item \textbf{Specjalista} -- zarządzanie własnym grafikiem, modyfikowanie dostępnych terminów, modyfikowanie grafiku także dla istniejących rezerwacji, przeglądanie rezerwacji dotyczących własnego grafiku.
  \item \textbf{Administrator} -- zarządzanie rolami użytkowników, konfiguracja zasad rezerwacji i anulowania, obsługa wyjątków od konfliktów rezerwacji, nadzór nad skalowalnością i poprawnym działaniem systemu.
\end{itemize}

% ─── WYMAGANIA FUNKCJONALNE ─────────────────────────────────────────
\section{Wymagania funkcjonalne}

\begin{description}[leftmargin=1.8cm, labelwidth=1.3cm, labelsep=0.5cm, style=nextline]
  \item[FR1] System musi umożliwiać użytkownikowi przeglądanie dostępnych terminów wizyt u specjalistów.
  \item[FR2] System musi umożliwiać użytkownikowi dokonanie rezerwacji wybranego dostępnego terminu.
  \item[FR3] System musi umożliwiać użytkownikowi anulowanie własnej rezerwacji najpóźniej do 24 godzin przed zaplanowanym terminem.
  \item[FR4] Po anulowaniu wizyty termin automatycznie wraca do puli dostępnych terminów.
  \item[FR5] System musi zapobiegać podwójnej rezerwacji: ten sam termin nie może zostać zarezerwowany przez więcej niż jednego użytkownika.
  \item[FR6] System musi umożliwiać obsługę wyjątków od standardowych reguł konfliktu rezerwacji w określonych przypadkach.
  \item[FR7] System musi umożliwiać specjaliście zarządzanie własnym grafikiem, w tym dodawanie, usuwanie i zmianę dostępnych terminów.
  \item[FR8] System musi umożliwiać specjaliście modyfikowanie grafiku także wtedy, gdy istnieją już powiązane rezerwacje.
  \item[FR9] Każda wizyta musi posiadać jeden ze zdefiniowanych statusów: AVAILABLE, BOOKED, CANCELLED, COMPLETED.
  \item[FR10] Jeden użytkownik może mieć maksymalnie 3 aktywne rezerwacje. Administrator może w uzasadnionych przypadkach zezwolić na przekroczenie tego limitu.
\end{description}

% ─── WYMAGANIA NIEFUNKCJONALNE ──────────────────────────────────────
\section{Wymagania niefunkcjonalne}

\begin{description}[leftmargin=1.8cm, labelwidth=1.3cm, labelsep=0.5cm, style=nextline]
  \item[NFR1] Operacja rezerwacji powinna zostać obsłużona w czasie krótszym niż 1 sekunda przy standardowym obciążeniu.
  \item[NFR2] System powinien zapewniać spójność danych i odporność na współbieżne próby rezerwacji tego samego terminu.
  \item[NFR3] System powinien skalować się wraz ze wzrostem liczby użytkowników, specjalistów, grafików i rezerwacji bez istotnego pogorszenia wydajności.
  \item[NFR4] Dostęp do operacji rezerwacji i anulowania wymaga uwierzytelnienia użytkownika.
  \item[NFR5] System powinien rejestrować operacje związane z rezerwacjami i anulowaniami w logu audytowym ze znacznikiem czasu w celu rozliczalności.
\end{description}

% ─── USER STORIES ───────────────────────────────────────────────────
\section{User stories}

\begin{description}[leftmargin=1.8cm, labelwidth=1.3cm, labelsep=0.5cm, style=nextline]
  \item[US1] Jako użytkownik chcę przeglądać dostępne terminy u specjalistów, aby wybrać dogodny termin wizyty.
  \item[US2] Jako użytkownik chcę zarezerwować dostępny termin wizyty, aby umówić spotkanie ze specjalistą.
  \item[US3] Jako użytkownik chcę anulować własną rezerwację na minimum 24h przed terminem, aby móc zrezygnować z wizyty.
  \item[US4] Jako specjalista chcę zarządzać własnym grafikiem, aby udostępniać i aktualizować terminy wizyt.
  \item[US5] Jako specjalista chcę modyfikować grafik także dla istniejących rezerwacji, aby dostosować harmonogram do zmian organizacyjnych.
  \item[US6] Jako administrator chcę konfigurować zasady rezerwacji i anulowania, aby dostosować działanie systemu do polityki organizacji.
  \item[US7] Jako administrator chcę obsługiwać wyjątki od konfliktów rezerwacji, aby umożliwiać niestandardowe przypadki zgodnie z obowiązującymi zasadami.
\end{description}

% ─── ACCEPTANCE CRITERIA ────────────────────────────────────────────
\section{Acceptance criteria}

\textbf{AC1}
\begin{itemize}[itemsep=2pt]
  \item \textbf{Given} użytkownik jest zalogowany i wybrany termin jest nadal dostępny do rezerwacji
  \item \textbf{When} użytkownik potwierdza rezerwację terminu wizyty
  \item \textbf{Then} system tworzy rezerwację i zmienia status terminu na BOOKED
\end{itemize}

\vspace{0.4em}
\textbf{AC2}
\begin{itemize}[itemsep=2pt]
  \item \textbf{Given} dwóch użytkowników jednocześnie próbuje zarezerwować ten sam termin
  \item \textbf{When} użytkownicy wysyłają żądanie rezerwacji w tym samym czasie
  \item \textbf{Then} system przyznaje rezerwację tylko jednemu użytkownikowi, zachowując spójność stanu końcowego
\end{itemize}

\vspace{0.4em}
\textbf{AC3}
\begin{itemize}[itemsep=2pt]
  \item \textbf{Given} użytkownik ma aktywną rezerwację, a do wizyty pozostało więcej niż 24 godziny
  \item \textbf{When} użytkownik anuluje rezerwację
  \item \textbf{Then} system zmienia status na CANCELLED oraz ponownie udostępnia termin
\end{itemize}

\vspace{0.4em}
\textbf{AC4}
\begin{itemize}[itemsep=2pt]
  \item \textbf{Given} użytkownik ma aktywną rezerwację, a do wizyty pozostało mniej niż 24 godziny
  \item \textbf{When} użytkownik próbuje anulować rezerwację
  \item \textbf{Then} system blokuje anulowanie i wyświetla informację o ograniczeniach
\end{itemize}

\vspace{0.4em}
\textbf{AC5}
\begin{itemize}[itemsep=2pt]
  \item \textbf{Given} użytkownik posiada już 3 aktywne rezerwacje
  \item \textbf{When} użytkownik próbuje zarezerwować kolejny termin
  \item \textbf{Then} system odrzuca operację, chyba że administrator zastosował wyjątek
\end{itemize}

\vspace{0.4em}
\textbf{AC6}
\begin{itemize}[itemsep=2pt]
  \item \textbf{Given} administrator jest zalogowany i posiada uprawnienia do konfiguracji zasad systemowych
  \item \textbf{When} administrator definiuje lub zmienia zasady rezerwacji oraz anulowania
  \item \textbf{Then} system zapisuje konfigurację i stosuje ją do nowych prób rezerwacji oraz anulowania
\end{itemize}

\vspace{0.4em}
\textbf{AC7}
\begin{itemize}[itemsep=2pt]
  \item \textbf{Given} występuje konflikt rezerwacji lub konflikt wynikający z reguł harmonogramu oraz administrator ma odpowiednie uprawnienia
  \item \textbf{When} administrator zatwierdza wyjątek zgodnie z obowiązującymi zasadami organizacji
  \item \textbf{Then} system pozwala na realizację operacji mimo standardowego konfliktu i zapisuje informację o zastosowanym wyjątku
\end{itemize}

\newpage

% ─── METRYKI VS GOLD ────────────────────────────────────────────────
\section{Metryki porównawcze względem gold artifact (Wpływ trybu Hybrid)}

\begin{mdframed}[backgroundcolor=metabg, linewidth=0pt, innerleftmargin=10pt, innerrightmargin=10pt, innertopmargin=8pt, innerbottommargin=8pt]
\small
\textbf{Wniosek z eksperymentu:} Interwencja manualna skutecznie wyeliminowała tzw. "gadulstwo" modelu (redukcja AC z 14 do 7, NFR z 8 do 5). Dzięki wstrzyknięciu do FR i AC logiki dla limitu 3 rezerwacji, 24-godzinnego okna oraz statusów wizyt, znacznie poprawiono metryki biznesowe artefaktu.
\end{mdframed}

\vspace{0.8em}

\subsection{Wymagania funkcjonalne (FR)}

\begin{tabular}{p{2.5cm} p{7cm} l}
\toprule
\textbf{ID Hybrid} & \textbf{Mapowanie na gold} & \textbf{Status} \\
\midrule
FR1 & Gold FR1 -- przeglądanie terminów & \tp{TP} \\
FR2 & Gold FR2 -- rezerwacja dostępnego terminu & \tp{TP} \\
FR3 & Gold FR5 -- anulowanie do 24h & \tp{TP} \\
FR4 & Gold FR6 -- zwolnienie terminu po anulowaniu & \tp{TP} \\
FR5 & Gold FR4 -- brak podwójnej rezerwacji & \tp{TP} \\
FR6 & brak w gold -- obsługa wyjątków w systemie & \fp{FP} \\
FR7 & Gold FR7 -- zarządzanie grafikiem specjalisty & \tp{TP} \\
FR8 & brak w gold -- modyfikacja istniejących rezerwacji & \fp{FP} \\
FR9 & Gold FR8 -- statusy wizyt (AVAILABLE, BOOKED...) & \tp{TP} \\
FR10 & Gold FR3 -- limit 3 aktywnych rezerwacji & \tp{TP} \\
\midrule
\multicolumn{3}{l}{\textit{Brakujące z gold:}} \\
-- & Gold FR9 -- dostęp do danych według ról & \fn{FN} \\
-- & Gold FR10 -- zakaz nakładających się rezerwacji & \fn{FN} \\
-- & Gold FR11 -- blokady slotów (BLOCKED) & \fn{FN} \\
-- & Gold FR13 -- spójność równoczesnych operacji & \fn{FN} \\
\bottomrule
\end{tabular}

\vspace{0.5em}
\begin{tabular}{lrrrr}
\toprule
& \textbf{TP} & \textbf{FP} & \textbf{FN} & \\
\midrule
FR & 8 & 2 & 4 & \\
\midrule
\textbf{Precision} & \multicolumn{4}{l}{$8 / (8+2) = \mathbf{0{,}80}$} \\
\textbf{Recall}    & \multicolumn{4}{l}{$8 / (8+4) = \mathbf{0{,}67}$} \\
\textbf{F1}        & \multicolumn{4}{l}{$2 \times 0{,}80 \times 0{,}67 / (0{,}80+0{,}67) = \mathbf{0{,}73}$} \\
\bottomrule
\end{tabular}

\vspace{1em}
\subsection{Wymagania niefunkcjonalne (NFR)}

\begin{tabular}{p{2.5cm} p{7cm} l}
\toprule
\textbf{ID Hybrid} & \textbf{Mapowanie na gold} & \textbf{Status} \\
\midrule
NFR1 (< 1 sek.) & Gold NFR2 -- wydajność & \tp{TP} \\
NFR2 (spójność) & Gold NFR1 -- spójność równoczesnych operacji & \tp{TP} \\
NFR3 (skalowanie)& Gold NFR5 -- skalowalność & \tp{TP} \\
NFR4 (uwierzyt.) & Gold NFR3 -- bezpieczeństwo (uwierzytelnienie) & \tp{TP} \\
NFR5 (log audytu)& Gold NFR4 -- log audytowy ze znacznikiem czasu & \tp{TP} \\
\midrule
\multicolumn{3}{l}{\textit{Brakujące z gold:}} \\
-- & Gold NFR6 -- zgodność z RODO / Szyfrowanie (ominięte) & \fn{FN} \\
\bottomrule
\end{tabular}

\vspace{0.5em}
\begin{tabular}{lrrrr}
\toprule
& \textbf{TP} & \textbf{FP} & \textbf{FN} & \\
\midrule
NFR & 5 & 0 & 1 & \\
\midrule
\textbf{Precision} & \multicolumn{4}{l}{$5 / (5+0) = \mathbf{1{,}00}$} \\
\textbf{Recall}    & \multicolumn{4}{l}{$5 / (5+1) = \mathbf{0{,}83}$} \\
\textbf{F1}        & \multicolumn{4}{l}{$2 \times 1{,}00 \times 0{,}83 / (1{,}00+0{,}83) = \mathbf{0{,}91}$} \\
\bottomrule
\end{tabular}

\vspace{1em}
\subsection{Acceptance criteria (AC)}

\begin{tabular}{p{2.5cm} p{7cm} l}
\toprule
\textbf{ID Hybrid} & \textbf{Mapowanie na gold} & \textbf{Status} \\
\midrule
AC1 & Gold AC1 -- rezerwacja $\rightarrow$ BOOKED & \tp{TP} \\
AC2 & Gold AC3 -- double booking spójność & \tp{TP} \\
AC3 & Gold AC4 -- anulowanie $>$24h $\rightarrow$ CANCELLED & \tp{TP} \\
AC4 & brak bezpośredniego (zablokowanie anulowania $<$24h) & \tp{TP} \\
AC5 & Gold AC2 -- limit 3 rezerwacji & \tp{TP} \\
AC6--AC7 & brak w gold -- scenariusze dla Administratora & \fp{FP} \\
\midrule
-- & Gold AC5 -- anulowanie graniczne (dokładnie 24h) & \fn{FN} \\
-- & Gold AC6 -- konflikt czasowy nakładających się rezerwacji & \fn{FN} \\
\bottomrule
\end{tabular}

\vspace{0.5em}
\begin{tabular}{lrrrr}
\toprule
& \textbf{TP} & \textbf{FP} & \textbf{FN} & \\
\midrule
AC & 5 & 2 & 2 & \\
\midrule
\textbf{Precision} & \multicolumn{4}{l}{$5 / (5+2) = \mathbf{0{,}71}$} \\
\textbf{Recall}    & \multicolumn{4}{l}{$5 / (5+2) = \mathbf{0{,}71}$} \\
\textbf{F1}        & \multicolumn{4}{l}{$2 \times 0{,}71 \times 0{,}71 / (0{,}71+0{,}71) = \mathbf{0{,}71}$} \\
\bottomrule
\end{tabular}

% ─── METADANE I OCENA ───────────────────────────────────────────────
\section{Metadane i ocena}
\begin{table}[h]
\begin{tabularx}{\textwidth}{l X}
\toprule
\textbf{Pole} & \textbf{Wartość} \\
\midrule
Model & \texttt{gpt-5.4} \\
Tryb & \textbf{Hybrid} (LLM generacja + Human-in-the-loop pruning) \\
Case study & appointment-booking (ver2 challenging) \\
\midrule
\multicolumn{2}{l}{\textit{Ocena (0--3)}} \\
\midrule
Correctness (M1)     & \scorebox{3} \quad błędy logiki biznesowej naprawione dzięki interwencji analityka \\
Completeness (M2)    & \scorebox{3} \quad dodano brakujące punkty węzłowe (np. czas do 24h, limit 3) \\
Consistency (M3)     & \scorebox{3} \quad artefakt jest w pełni spójny \\
Clarity (M4)         & \scorebox{3} \quad doskonała czytelność; radykalnie ograniczono objętość dokumentu \\
Maintainability (M5) & \scorebox{3} \quad wysoce użyteczny dokument gotowy do analizy architektonicznej \\
\bottomrule
\end{tabularx}
\end{table}

\end{document}